\section{Terminology}
\label{sec:terminology}
\slidesrc{Terminology introduction: 34--51; Semantics (long version): 44--50}

Every field has its vocabulary, and NLP borrows a lot of it from linguistics (morpheme,
corpus, collocation) and from information retrieval (precision, recall). These terms
are used throughout the course without being redefined, so it pays to fix them now.

This section covers:
\begin{enumerate}[nosep]
  \item the units words are made of: morphemes, words and compounds
        (Section~\ref{sec:morphology});
  \item a recap of the parts of speech (noun, verb, adjective, \ldots) with their Italian
        names and Penn tags (Section~\ref{sec:parts-of-speech});
  \item collections of texts and what we observe in them: corpora, genre, collocations
        (Section~\ref{sec:corpora});
  \item how to measure the quality of a system that retrieves documents: precision and
        recall (Section~\ref{sec:precision-recall}).
\end{enumerate}

\subsection{Morphemes, Words and Compounds}
\label{sec:morphology}
\slidesrc{Terminology introduction: 35--38, 42}

\subsubsection{Morphemes}
\slidesrc{Terminology introduction: 35--36}

\begin{nbdef}{Morpheme}
A \textbf{morpheme} is the smallest grammatical unit in a language; in other words, the
smallest \emph{meaningful} unit of a language. The field of linguistics that studies
morphemes is called \textbf{morphology}.
\end{nbdef}

A morpheme is not the same as a word. The principal difference is that a morpheme may
or may not stand alone, whereas a word, by definition, is freestanding. In
\emph{dogs} there is one word but two morphemes: \emph{dog} (the animal) and \emph{-s}
(``more than one''). Neither \emph{d} nor \emph{og} carries meaning, so they are not
morphemes.

\begin{nbdef}{Free and bound morphemes}
Every morpheme is either free or bound.
\begin{itemize}[nosep]
  \item A \textbf{free morpheme} can function independently as a word (e.g.\
        \emph{town}, \emph{dog}).
  \item A \textbf{bound morpheme} appears only as part of a word, always together with a
        \textbf{root} and sometimes with other bound morphemes. For example,
        \emph{un-} never appears alone. Most bound morphemes in English are
        \textbf{affixes}: \textbf{prefixes} (before the root, \emph{un-}) and
        \textbf{suffixes} (after the root, \emph{-ness}).
\end{itemize}
\end{nbdef}

\begin{nbexample}{unbreakable}
\emph{Unbreakable} is made of three morphemes:
\begin{center}
\begin{tikzpicture}[
    m/.style={draw, rounded corners, minimum height=0.8cm, minimum width=1.7cm, font=\small},
    lab/.style={font=\scriptsize, align=center}]
  \node[font=\large\itshape] (w) at (2.4,1.5) {unbreakable};
  \node[m, fill=nbred!8,  draw=nbred]  (a) at (0,0)   {un-};
  \node[m, fill=nbblue!8, draw=nbblue] (b) at (2.4,0) {break};
  \node[m, fill=nbred!8,  draw=nbred]  (c) at (4.8,0) {-able};
  \draw[nbgray] (w.south) -- (a.north); \draw[nbgray] (w.south) -- (b.north); \draw[nbgray] (w.south) -- (c.north);
  \node[lab, below=2pt of a] {bound prefix\\``not''};
  \node[lab, below=2pt of b] {root\\free morpheme};
  \node[lab, below=2pt of c] {bound suffix\\``can be done''};
\end{tikzpicture}
\end{center}
\end{nbexample}

\begin{nbwarning}[\emph{-able} is bound]
The slide calls \emph{-able} ``a free morpheme signifying \emph{can be done}''. As a
suffix it is \textbf{bound}: it cannot stand alone with the meaning ``that can be
done'' and needs a root (\emph{break-able}, \emph{read-able}). The English word
\emph{able} exists, but it is a separate free morpheme with a different use (``she is
able to swim''). So \emph{unbreakable} = bound prefix + free root + bound suffix.
\end{nbwarning}

\subsubsection{Derivational and inflectional morphemes}
\slidesrc{Terminology introduction: 37}

Bound morphemes are further divided according to \emph{what they change}.

\begin{nbdef}{Derivational and inflectional morphemes}
\begin{itemize}[nosep]
  \item A \textbf{derivational morpheme}, combined with a root, changes the
        \textbf{meaning} or the \textbf{part of speech} of the word: it creates a
        \emph{new word}.
  \item An \textbf{inflectional morpheme} modifies a verb's tense, aspect, mood, person or
        number, or a noun's, pronoun's or adjective's number, gender or case,
        \emph{without} changing the word's meaning or class (part of speech): it creates
        a different \emph{form of the same word}.
\end{itemize}
\end{nbdef}

\begin{center}
\small
\begin{tabularx}{\linewidth}{@{}l l l X@{}}
\toprule
\textbf{Root} & \textbf{Morpheme} & \textbf{Result} & \textbf{What changed} \\
\midrule
\multicolumn{4}{@{}l}{\emph{Derivational}} \\
happy (adjective) & -ness & happiness (noun) & part of speech: adjective $\to$ noun \\
kind              & un-   & unkind           & meaning is inverted \\
\midrule
\multicolumn{4}{@{}l}{\emph{Inflectional}} \\
dog  & -s  & dogs   & number (singular $\to$ plural) \\
wait & -ed & waited & tense (present $\to$ past) \\
\bottomrule
\end{tabularx}
\captionof{table}{Derivation makes a new word, inflection a new form of the same word.}
\end{center}

\begin{nbintuition}[a dictionary test]
Ask: ``would a dictionary give the result its own entry?'' \emph{Happiness} and
\emph{unkind} have their own entries (derivation), \emph{dogs} and \emph{waited} are found
under \emph{dog} and \emph{wait} (inflection). This is exactly the difference between
\emph{lemmatisation} (going back to the dictionary form, removing inflection) and the
analysis of word formation, which we will meet later in the course.
\end{nbintuition}

\subsubsection{Words}
\slidesrc{Terminology introduction: 38}

\begin{nbdef}{Word}
A \textbf{word} is the smallest element that can be uttered in isolation with
objective or practical meaning.
\end{nbdef}

This contrasts with a morpheme, which is the smallest unit of meaning but does not
necessarily stand on its own:
\begin{itemize}[nosep]
  \item a word may consist of a \textbf{single morpheme}: \emph{oh!}, \emph{rock},
        \emph{red}, \emph{quick}, \emph{run}, \emph{expect};
  \item or of \textbf{several}: \emph{rock-s}, \emph{red-ness}, \emph{quick-ly},
        \emph{runn-ing}, \emph{un-expect-ed};
  \item and a morpheme may be unable to stand on its own as a word: in the examples,
        \emph{-s}, \emph{-ness}, \emph{-ly}, \emph{-ing}, \emph{un-}, \emph{-ed}.
\end{itemize}

\subsubsection{Compounds}
\slidesrc{Terminology introduction: 42}

\begin{nbdef}{Compound}
A \textbf{compound} is a word composed of more than one \emph{free} morpheme.
\end{nbdef}

The last part is usually the \textbf{head} (it gives the category of the whole word),
and the first part is the \textbf{modifier}. Compounds can be written as one word
(\emph{football}), with a hyphen (\emph{blue-green}) or even as two words
(\emph{snow white}).

\begin{center}
\small
\begin{tabular}{@{}l l l @{\hspace{1.6cm}} l l l@{}}
\toprule
\textbf{Modifier} & \textbf{Head} & \textbf{Compound} & \textbf{Modifier} & \textbf{Head} & \textbf{Compound} \\
\midrule
noun        & noun      & football    & noun        & verb        & browbeat \\
adjective   & noun      & blackboard  & adjective   & verb        & highlight \\
verb        & noun      & breakwater  & verb        & verb        & freeze-dry \\
preposition & noun      & underworld  & preposition & verb        & undercut \\
noun        & adjective & snow white  & noun        & preposition & love-in \\
adjective   & adjective & blue-green  & adverb      & preposition & forthwith \\
verb        & adjective & tumbledown  & verb        & preposition & takeout \\
preposition & adjective & over-ripe   & preposition & preposition & without \\
\bottomrule
\end{tabular}
\captionof{table}{English compounds by the category of their two parts.}
\end{center}

\begin{nbwarning}[compound vs.\ derived word]
\emph{Blackboard} is a compound (\emph{black} and \emph{board} are both free).
\emph{Unkind} is \emph{not} a compound: \emph{un-} is bound, so it is a derived word.
Also note that the meaning of a compound is often not the sum of its parts: a
\emph{blackboard} can be green.
\end{nbwarning}

\subsection{Parts of Speech: a Grammar Recap}
\label{sec:parts-of-speech}
\notinslides

The course often talks about nouns, pronouns, adjectives, determiners\ldots{} without
defining them: they are assumed from school grammar. This subsection is a quick recap to
come back to whenever needed, with the link to the Italian grammar you already know and
the Penn Treebank tags used in POS tagging (Section~\ref{sec:pos-tagging}). Italian
translations are always in \ita{green}; most tables pair each English word with its
Italian counterpart, one per row.

\begin{nbintuition}[English and Italian names]
The English \textbf{parts of speech} (also \emph{word classes} or \emph{lexical
categories}) are the Italian \textbf{parti del discorso}. Assigning one to each word of a
sentence is what Italian school calls \textbf{analisi grammaticale}, and what NLP calls
\textbf{POS tagging}. The sentence-level analysis (subject, object, \ldots) is the
\textbf{analisi logica}, recalled at the end of this subsection.
\end{nbintuition}

\subsubsection{Overview}
\notinslides

\begin{center}
\footnotesize
\begin{tabularx}{\linewidth}{@{}l >{\raggedright\arraybackslash}p{2.5cm} c >{\raggedright\arraybackslash}p{2.6cm} >{\raggedright\arraybackslash}X@{}}
\toprule
\textbf{English} & \textbf{Italiano} & \textbf{Class} & \textbf{Penn tags} & \textbf{Examples} \\
\midrule
Noun         & nome (sostantivo) & open   & NN, NNS, NNP, NNPS & cat \ita{gatto}, ideas \ita{idee}, Rome \ita{Roma} \\
Verb         & verbo             & open   & VB, VBD, VBG, VBN, VBP, VBZ, MD & run \ita{correre}, ran \ita{corse}, is \ita{è}, can \ita{può} \\
Adjective    & aggettivo (qualificativo) & open & JJ, JJR, JJS & red \ita{rosso}, bigger \ita{più grande}, best \ita{il migliore} \\
Adverb       & avverbio          & open   & RB, RBR, RBS, WRB & quickly \ita{velocemente}, very \ita{molto}, where \ita{dove} \\
Pronoun      & pronome           & closed & PRP, PRP\$, WP, WP\$, EX & I \ita{io}, me \ita{me/mi}, mine \ita{il mio}, who \ita{chi} \\
Determiner   & articolo, agg.\ determinativo & closed & DT, PDT, WDT & the \ita{il/la}, a \ita{un/una}, this \ita{questo}, every \ita{ogni} \\
Preposition  & preposizione      & closed & IN, TO & of \ita{di}, to \ita{a}, from \ita{da}, in \ita{in} \\
Conjunction  & congiunzione      & closed & CC, IN & and \ita{e}, but \ita{ma}, because \ita{perché} \\
Interjection & interiezione      & closed & UH & oh \ita{oh}, wow \ita{caspita}, ouch \ita{ahi} \\
\midrule
Numeral      & aggettivo numerale & closed & CD & three \ita{tre}, 1908 \\
Particle     & (no equivalent)   & closed & RP & give \emph{up} \ita{rinunciare}, turn \emph{off} \ita{spegnere} \\
\bottomrule
\end{tabularx}
\captionof{table}{Parts of speech at a glance. \emph{Open} classes accept new words all
the time (\emph{to google} \ita{googlare}, \emph{selfie}); \emph{closed} classes are small
fixed lists (Section~\ref{sec:grammars-parsing}).}
\end{center}

\begin{nbwarning}[variable and invariable]
In Italian, five parts of speech are \emph{variabili} (articolo, nome, aggettivo,
pronome, verbo: they change with gender, number, person, tense) and four are
\emph{invariabili} (avverbio, preposizione, congiunzione, interiezione). English inflects
much less: \textbf{articles and adjectives never change} (\emph{the red car}
\ita{la macchina rossa}, \emph{the red cars} \ita{le macchine rosse}: in English only
\emph{car} changes), and nouns have number but \textbf{no grammatical gender}.
\end{nbwarning}

%======================================================================
\subsubsection{Nouns}
\notinslides

A \textbf{noun} names a person, animal, thing, place or idea. \quad
\textbf{Italiano}: \emph{nome} / \emph{sostantivo}.

\begin{center}
\footnotesize
\begin{tabularx}{\linewidth}{@{}>{\raggedright\arraybackslash}p{3.3cm} >{\raggedright\arraybackslash}p{3.2cm} >{\raggedright\arraybackslash}X l@{}}
\toprule
\textbf{Type} & \textbf{Italiano} & \textbf{Examples} & \textbf{Penn} \\
\midrule
common             & nome comune            & city \ita{città}, dog \ita{cane} & NN \\
proper             & nome proprio           & Genoa \ita{Genova}, Mary \ita{Maria} & NNP \\
plural             & plurale                & dogs \ita{cani}, children \ita{bambini}, the Alps \ita{le Alpi} & NNS, NNPS \\
countable          & numerabile             & a chair \ita{una sedia}, two chairs \ita{due sedie} & NN, NNS \\
uncountable (mass) & non numerabile         & water \ita{acqua}, information \ita{informazioni}, advice \ita{consigli} & NN \\
concrete / abstract & concreto / astratto   & table \ita{tavolo} / happiness \ita{felicità} & NN \\
compound           & nome composto          & football \ita{calcio}, bus stop \ita{fermata dell'autobus} & NN \\
possessive \emph{'s} & complemento di specificazione & Mary's book \ita{il libro di Maria} & NNP + POS \\
\bottomrule
\end{tabularx}
\end{center}
\textbf{Differences with Italian}: no gender (\emph{the table} \ita{il tavolo},
\emph{the chair} \ita{la sedia}: the same \emph{the}); uncountable nouns have no plural and
no \emph{a} even when Italian uses a plural (\emph{some information} \ita{delle
informazioni}, \emph{a piece of advice} \ita{un consiglio}).

%======================================================================
\subsubsection{Pronouns}
\notinslides

A \textbf{pronoun} stands for a noun (or a whole noun phrase). \quad
\textbf{Italiano}: \emph{pronome}.

\needspace{8\baselineskip}\paragraph{Personal, possessive and reflexive pronouns.} One row per person:

\begin{center}
\scriptsize
\setlength{\tabcolsep}{3pt}
\begin{tabular}{@{}l l l l l l@{}}
\toprule
\textbf{Person} & \textbf{Subject} & \textbf{Object} & \textbf{Possessive adj.} & \textbf{Possessive pron.} & \textbf{Reflexive} \\
 & \emph{pers.\ soggetto} & \emph{pers.\ complemento} & \emph{agg.\ possessivo} & \emph{pron.\ possessivo} & \emph{riflessivo} \\
\midrule
1s & I \ita{io}      & me \ita{me, mi}       & my \ita{mio}     & mine \ita{il mio}    & myself \ita{mi, me stesso} \\
2s & you \ita{tu}    & you \ita{te, ti}      & your \ita{tuo}   & yours \ita{il tuo}   & yourself \ita{ti, te stesso} \\
3s m. & he \ita{lui}    & him \ita{lui, lo, gli} & his \ita{suo}   & his \ita{il suo}     & himself \ita{si, sé stesso} \\
3s f. & she \ita{lei}   & her \ita{lei, la, le} & her \ita{suo}    & hers \ita{il suo}    & herself \ita{si, sé stessa} \\
3s thing & it \ita{esso, --} & it \ita{lo, la}     & its \ita{suo}    & --                   & itself \ita{si} \\
1p & we \ita{noi}    & us \ita{noi, ci}      & our \ita{nostro} & ours \ita{il nostro} & ourselves \ita{ci, noi stessi} \\
2p & you \ita{voi}   & you \ita{voi, vi}     & your \ita{vostro} & yours \ita{il vostro} & yourselves \ita{vi, voi stessi} \\
3p & they \ita{loro} & them \ita{loro, li, le} & their \ita{loro} & theirs \ita{il loro} & themselves \ita{si, sé stessi} \\
\midrule
\textbf{Penn tag} & PRP & PRP & PRP\$ & PRP & PRP \\
\bottomrule
\end{tabular}
\end{center}

\needspace{8\baselineskip}\paragraph{The other pronouns.}
\begin{center}
\footnotesize
\begin{tabularx}{\linewidth}{@{}l >{\raggedright\arraybackslash}p{3.0cm} >{\raggedright\arraybackslash}X l@{}}
\toprule
\textbf{Type} & \textbf{English} & \textbf{Italiano} & \textbf{Penn} \\
\midrule
Demonstrative & this / these           & \ita{questo, questa / questi, queste} & DT \\
              & that / those           & \ita{quello, quella / quelli, quelle} & DT \\
\midrule
Relative      & who (people)           & \ita{che, il quale} \quad \emph{the man \underline{who} called} \ita{l'uomo che ha chiamato} & WP \\
              & which (things)         & \ita{che, il quale} \quad \emph{the book \underline{which} I read} \ita{il libro che ho letto} & WDT \\
              & that (people, things)  & \ita{che} & WDT \\
              & whom                   & \ita{cui} \quad \emph{the girl \underline{to whom} I wrote} \ita{la ragazza a cui ho scritto} & WP \\
              & whose                  & \ita{il cui} \quad \emph{the man \underline{whose} car\ldots} \ita{l'uomo la cui macchina\ldots} & WP\$ \\
\midrule
Interrogative & who?                   & \ita{chi?} & WP \\
              & what?                  & \ita{che cosa?} & WP \\
              & which?                 & \ita{quale?} & WDT \\
              & whose?                 & \ita{di chi?} & WP\$ \\
\midrule
Indefinite    & someone / somebody     & \ita{qualcuno} & NN \\
              & something              & \ita{qualcosa} & NN \\
              & nobody / no one        & \ita{nessuno} & NN \\
              & nothing                & \ita{niente, nulla} & NN \\
              & everybody / everyone   & \ita{tutti, ognuno} & NN \\
              & everything             & \ita{tutto} & NN \\
              & anybody / anything     & \ita{chiunque / qualsiasi cosa} (in questions: \ita{qualcuno / qualcosa}) & NN \\
\bottomrule
\end{tabularx}
\end{center}

\textbf{Differences with Italian}:
\begin{itemize}[nosep]
  \item Italian is \emph{pro-drop}: the subject pronoun can be omitted (\ita{mangio}).
        English always needs it (\emph{I eat}), even a dummy one: \emph{\underline{it}
        rains} \ita{piove}, \emph{\underline{there} is a cat} \ita{c'è un gatto}
        (\emph{there} here is ``existential'', tag EX).
  \item The possessive adjective has \emph{no article} in English: \emph{my book}
        \ita{il mio libro}; and it agrees with the \emph{owner}: \emph{her brother}
        \ita{suo fratello} (of a woman), \emph{his sister} \ita{sua sorella} (of a man).
  \item Italian unstressed object pronouns (\ita{mi, ti, lo, la, ci, vi, li, le, gli})
        go before the verb, English ones after it: \ita{lo vedo} \emph{I see him}.
\end{itemize}

%======================================================================
\subsubsection{Verbs}
\notinslides

A \textbf{verb} expresses an action, an event or a state. \quad
\textbf{Italiano}: \emph{verbo}.

\needspace{8\baselineskip}\paragraph{The forms of an English verb.} English has only a few verb forms, and they
are exactly what the Penn tags list:
\begin{center}
\footnotesize
\begin{tabularx}{\linewidth}{@{}l >{\raggedright\arraybackslash}p{2.4cm} >{\raggedright\arraybackslash}X >{\raggedright\arraybackslash}p{3.2cm}@{}}
\toprule
\textbf{Tag} & \textbf{Form} & \textbf{Example} & \textbf{Italian equivalent} \\
\midrule
VB  & base form             & \emph{to \underline{eat}} \ita{mangiare}, \emph{I can \underline{eat}} \ita{posso mangiare} & infinito \\
VBP & present, not 3rd sing. & \emph{I/we/they \underline{eat}} \ita{mangio / mangiamo / mangiano} & presente indicativo \\
VBZ & present, 3rd sing.    & \emph{she \underline{eats}} \ita{(lei) mangia} & presente, 3\textsuperscript{a} sing. \\
VBD & past tense            & \emph{she \underline{ate}} \ita{mangiò / mangiava / ha mangiato} & passato remoto / imperfetto / pass.\ prossimo \\
VBN & past participle       & \emph{she has \underline{eaten}} \ita{ha mangiato} & participio passato \\
VBG & gerund / present participle & \emph{she is \underline{eating}} \ita{sta mangiando}, \emph{\underline{eating} is fun} \ita{mangiare è divertente} & gerundio, infinito sostantivato \\
\bottomrule
\end{tabularx}
\end{center}

\needspace{8\baselineskip}\paragraph{English tenses and Italian tempi.} Everything beyond those forms is built
with auxiliaries (\emph{be, have, do}) and modals:
\begin{center}
\footnotesize
\begin{tabularx}{\linewidth}{@{}>{\raggedright\arraybackslash}p{3.0cm} l >{\raggedright\arraybackslash}X >{\raggedright\arraybackslash}p{3.0cm}@{}}
\toprule
\textbf{English tense} & \textbf{Example} & \textbf{Italiano} & \textbf{Tempo / modo} \\
\midrule
present simple       & I eat              & \ita{mangio}                    & presente \\
present continuous   & I am eating        & \ita{sto mangiando / mangio}    & presente (stare + gerundio) \\
past simple          & I ate              & \ita{mangiai / ho mangiato}     & passato remoto / prossimo \\
past continuous      & I was eating       & \ita{mangiavo / stavo mangiando} & imperfetto \\
present perfect      & I have eaten       & \ita{ho mangiato}               & passato prossimo \\
past perfect         & I had eaten        & \ita{avevo mangiato}            & trapassato prossimo \\
future               & I will eat         & \ita{mangerò}                   & futuro semplice \\
future perfect       & I will have eaten  & \ita{avrò mangiato}             & futuro anteriore \\
conditional          & I would eat        & \ita{mangerei}                  & condizionale presente \\
past conditional     & I would have eaten & \ita{avrei mangiato}            & condizionale passato \\
imperative           & eat!               & \ita{mangia!}                   & imperativo \\
passive              & it is / was eaten  & \ita{è / fu mangiato}           & forma passiva \\
\bottomrule
\end{tabularx}
\end{center}
The Italian \emph{congiuntivo} has (almost) no English form: \emph{I think he is right}
\ita{penso che abbia ragione}, \emph{if I were rich} \ita{se fossi ricco}.

\needspace{8\baselineskip}\paragraph{Auxiliaries and modals.}
\begin{center}
\footnotesize
\begin{tabularx}{\linewidth}{@{}l >{\raggedright\arraybackslash}p{3.6cm} >{\raggedright\arraybackslash}X l@{}}
\toprule
\textbf{English} & \textbf{Italiano} & \textbf{Example} & \textbf{Penn} \\
\midrule
be      & \ita{essere} (also \ita{stare})  & \emph{she is tired} \ita{è stanca}, \emph{she is eating} \ita{sta mangiando} & VB\ldots \\
have    & \ita{avere}                      & \emph{she has eaten} \ita{ha mangiato} & VB\ldots \\
do      & -- (no equivalent)               & \emph{do you know?} \ita{sai?}, \emph{I don't know} \ita{non so} & VB\ldots \\
can     & \ita{potere, sapere}             & \emph{I can swim} \ita{so nuotare}, \emph{can I go?} \ita{posso andare?} & MD \\
must    & \ita{dovere}                     & \emph{you must go} \ita{devi andare} & MD \\
should  & \ita{dovrei} (condizionale)      & \emph{you should rest} \ita{dovresti riposare} & MD \\
may / might & \ita{potrei, può darsi che}  & \emph{it might rain} \ita{potrebbe piovere} & MD \\
will    & (futuro)                         & \emph{she will come} \ita{verrà} & MD \\
would   & (condizionale)                   & \emph{she would come} \ita{verrebbe} & MD \\
\bottomrule
\end{tabularx}
\end{center}
Modals never change form (\emph{she can}, not \emph{she cans}) and are followed by the
base form (VB). They correspond to the Italian \emph{verbi servili} (\ita{potere, dovere,
volere}) or to a tense of the main verb.

\needspace{8\baselineskip}\paragraph{Other verb distinctions.}
\begin{itemize}[nosep]
  \item \textbf{Transitive} \ita{transitivo}: \emph{she reads \underline{a book}}
        \ita{lei legge un libro}; \textbf{intransitive} \ita{intransitivo}: \emph{she
        sleeps} \ita{lei dorme}.
  \item \textbf{Phrasal verbs}: verb + particle (RP) with a new meaning: \emph{give up}
        \ita{rinunciare}, \emph{look after} \ita{badare a}, \emph{turn off}
        \ita{spegnere}, \emph{find out} \ita{scoprire}.
\end{itemize}

%======================================================================
\subsubsection{Adjectives}
\notinslides

An \textbf{adjective} describes or qualifies a noun. \quad
\textbf{Italiano}: \emph{aggettivo qualificativo}.

\begin{center}
\footnotesize
\begin{tabularx}{\linewidth}{@{}>{\raggedright\arraybackslash}p{3.6cm} >{\raggedright\arraybackslash}X >{\raggedright\arraybackslash}X l@{}}
\toprule
\textbf{Degree} \emph{(grado)} & \textbf{Short adjective} & \textbf{Long adjective} & \textbf{Penn} \\
\midrule
positive \emph{(positivo)}      & big \ita{grande}                  & expensive \ita{costoso}                   & JJ \\
comparative of majority \emph{(comp.\ di maggioranza)} & bigger than \ita{più grande di} & more expensive than \ita{più costoso di} & JJR / RBR + JJ \\
comparative of minority \emph{(comp.\ di minoranza)}   & less big than \ita{meno grande di} & less expensive than \ita{meno costoso di} & RBR + JJ \\
comparative of equality \emph{(comp.\ di uguaglianza)} & as big as \ita{grande come}   & as expensive as \ita{costoso come}        & RB + JJ \\
relative superlative \emph{(superlativo relativo)}     & the biggest \ita{il più grande} & the most expensive \ita{il più costoso} & JJS / RBS + JJ \\
absolute superlative \emph{(superlativo assoluto)}     & very big \ita{grandissimo}    & very expensive \ita{costosissimo}         & RB + JJ \\
\bottomrule
\end{tabularx}
\end{center}

Irregular: \emph{good, better, best} \ita{buono, migliore, il migliore};
\emph{bad, worse, worst} \ita{cattivo, peggiore, il peggiore}; \emph{far, further,
furthest} \ita{lontano, più lontano, il più lontano}.

Position: \textbf{attributive} \emph{a \underline{red} car} \ita{una macchina rossa}
vs.\ \textbf{predicative} \emph{the car is \underline{red}} \ita{la macchina è rossa}
(\emph{funzione attributiva / predicativa}).

\textbf{Differences with Italian}: English adjectives are invariable and come
\emph{before} the noun: \emph{a red car, two red cars} \ita{una macchina rossa, due
macchine rosse}. What Italian calls \emph{aggettivi determinativi} (possessivi,
dimostrativi, indefiniti, interrogativi, numerali) are \textbf{not} adjectives in English
grammar: \emph{my} \ita{mio} is PRP\$, \emph{this} \ita{questo} and \emph{every}
\ita{ogni} are DT, \emph{three} \ita{tre} is CD.

%======================================================================
\subsubsection{Adverbs}
\notinslides

An \textbf{adverb} modifies a verb, an adjective, another adverb or a whole sentence.
\quad \textbf{Italiano}: \emph{avverbio}. Manner adverbs are formed with \emph{-ly}, the
Italian \emph{-mente}: \emph{quick-ly} \ita{veloce-mente}.

\begin{center}
\footnotesize
\begin{tabularx}{\linewidth}{@{}>{\raggedright\arraybackslash}p{3.2cm} >{\raggedright\arraybackslash}X l@{}}
\toprule
\textbf{Type} & \textbf{English $\to$ Italiano} & \textbf{Penn} \\
\midrule
manner \emph{(di modo)}           & quickly \ita{velocemente}, slowly \ita{lentamente}, well \ita{bene}, badly \ita{male} & RB \\
time \emph{(di tempo)}            & now \ita{ora}, today \ita{oggi}, yesterday \ita{ieri}, tomorrow \ita{domani}, soon \ita{presto}, already \ita{già}, still \ita{ancora} & RB, NN \\
place \emph{(di luogo)}           & here \ita{qui}, there \ita{lì}, everywhere \ita{ovunque}, inside \ita{dentro}, outside \ita{fuori} & RB \\
frequency \emph{(di frequenza)}   & always \ita{sempre}, often \ita{spesso}, sometimes \ita{a volte}, never \ita{mai} & RB \\
quantity \emph{(di quantità)}     & very \ita{molto}, too \ita{troppo}, quite \ita{abbastanza}, a little \ita{un po'}, almost \ita{quasi} & RB \\
affirmation / negation            & yes \ita{sì}, not \ita{non}, also \ita{anche}, certainly \ita{certamente} & RB, UH \\
interrogative \emph{(interrogativo)} & when? \ita{quando?}, where? \ita{dove?}, why? \ita{perché?}, how? \ita{come?}, how much? \ita{quanto?} & WRB \\
\bottomrule
\end{tabularx}
\end{center}
Comparison works as for adjectives: \emph{fast} \ita{velocemente}, \emph{faster} \ita{più
velocemente} (RBR), \emph{fastest} \ita{il più velocemente} (RBS).

%======================================================================
\subsubsection{Determiners and articles}
\notinslides

A \textbf{determiner} comes before a noun and specifies \emph{which} or \emph{how many}.
\quad \textbf{Italiano}: \emph{articolo} and most \emph{aggettivi determinativi}.

\begin{center}
\footnotesize
\begin{tabularx}{\linewidth}{@{}>{\raggedright\arraybackslash}p{3.0cm} l >{\raggedright\arraybackslash}X l@{}}
\toprule
\textbf{Type} & \textbf{English} & \textbf{Italiano} & \textbf{Penn} \\
\midrule
definite article \emph{(determinativo)}  & the     & \ita{il, lo, la, l', i, gli, le}: \emph{the boy, the girl, the books} \ita{il ragazzo, la ragazza, i libri} & DT \\
indefinite article \emph{(indeterminativo)} & a / an & \ita{un, uno, una, un'}: \emph{a dog, an apple, an hour} \ita{un cane, una mela, un'ora} & DT \\
partitive \emph{(partitivo)}    & some / --   & \ita{del, dello, della, dei, degli, delle}: \emph{some bread} \ita{del pane} & DT \\
\midrule
demonstrative \emph{(dimostrativo)} & this / these & \ita{questo, questa / questi, queste} & DT \\
                                    & that / those & \ita{quel, quello, quella / quei, quegli, quelle} & DT \\
\midrule
quantifier \emph{(indefinito)}  & every / each & \ita{ogni / ciascuno}: \emph{every day} \ita{ogni giorno} & DT \\
                                & some / any   & \ita{qualche, alcuni / qualsiasi} & DT \\
                                & no           & \ita{nessuno}: \emph{no money} \ita{nessun soldo} & DT \\
                                & all / both   & \ita{tutti / entrambi} & DT, PDT \\
                                & many / much  & \ita{molti / molto}: \emph{many books, much time} \ita{molti libri, molto tempo} & JJ \\
                                & few / little & \ita{pochi / poco} & JJ \\
\midrule
interrogative \emph{(interrogativo)} & which? / what? & \ita{quale? / che?}: \emph{which book?} \ita{quale libro?} & WDT \\
                                     & whose?          & \ita{di chi?}: \emph{whose car?} \ita{di chi è la macchina?} & WP\$ \\
\bottomrule
\end{tabularx}
\end{center}

The English articles never change: one \emph{the} for all seven Italian forms, and the
choice between \emph{a} and \emph{an} depends only on the \emph{sound} that follows
(\emph{an hour}, \emph{a university}). Predeterminers (PDT) come before another
determiner: \emph{\underline{all} the cats} \ita{tutti i gatti}, \emph{\underline{both}
the girls} \ita{entrambe le ragazze}.

%======================================================================
\subsubsection{Prepositions}
\notinslides

A \textbf{preposition} links a noun phrase to the rest of the sentence, expressing place,
time, means, \ldots \quad \textbf{Italiano}: \emph{preposizione}.

\needspace{8\baselineskip}\paragraph{The Italian preposizioni semplici and their English translations.} There is
no one-to-one mapping: each Italian preposition has several English translations,
depending on the meaning.

\begin{center}
\footnotesize
\begin{tabularx}{\linewidth}{@{}l l >{\raggedright\arraybackslash}X@{}}
\toprule
\textbf{Italiano} & \textbf{English} & \textbf{Examples} \\
\midrule
\ita{di}      & of          & \emph{a glass of wine} \ita{un bicchiere di vino} \\
              & from        & \emph{I am from Genoa} \ita{sono di Genova} \\
              & about       & \emph{let's talk about it} \ita{parliamone (parliamo di questo)} \\
              & than        & \emph{taller than me} \ita{più alto di me} \\
              & 's          & \emph{Mary's book} \ita{il libro di Maria} \\
\midrule
\ita{a}       & to          & \emph{I go to Rome} \ita{vado a Roma}, \emph{I give it to Bob} \ita{lo do a Bob} \\
              & at          & \emph{at home} \ita{a casa}, \emph{at 5} \ita{alle 5} \\
              & in          & \emph{I live in Genoa} \ita{abito a Genova} \\
\midrule
\ita{da}      & from        & \emph{from Rome to Milan} \ita{da Roma a Milano} \\
              & by          & \emph{written by Dante} \ita{scritto da Dante} \\
              & for / since & \emph{for two years} \ita{da due anni}, \emph{since 2020} \ita{dal 2020} \\
              & at / to someone's & \emph{at Mario's} \ita{da Mario}, \emph{to the doctor} \ita{dal dottore} \\
\midrule
\ita{in}      & in / into   & \emph{in Italy} \ita{in Italia}, \emph{into the room} \ita{nella stanza} \\
              & by          & \emph{by train} \ita{in treno} \\
\midrule
\ita{con}     & with        & \emph{with Bob} \ita{con Bob}, \emph{with a knife} \ita{con un coltello} \\
\midrule
\ita{su}      & on          & \emph{on the table} \ita{sul tavolo} \\
              & about       & \emph{a book about Rome} \ita{un libro su Roma} \\
              & out of      & \emph{one out of ten} \ita{uno su dieci} \\
\midrule
\ita{per}     & for         & \emph{for you} \ita{per te}, \emph{for two hours} \ita{per due ore} \\
              & through     & \emph{through the city} \ita{per la città} \\
              & to (purpose) & \emph{I study to learn} \ita{studio per imparare} \\
\midrule
\ita{tra / fra} & between   & \emph{between Rome and Milan} \ita{tra Roma e Milano} \\
              & among       & \emph{among friends} \ita{tra amici} \\
              & in (time)   & \emph{in two hours} \ita{tra due ore} \\
\bottomrule
\end{tabularx}
\end{center}

\needspace{8\baselineskip}\paragraph{Other common English prepositions.} \emph{under} \ita{sotto}, \emph{over /
above} \ita{sopra}, \emph{near} \ita{vicino a}, \emph{behind} \ita{dietro}, \emph{in front
of} \ita{davanti a}, \emph{before} \ita{prima di}, \emph{after} \ita{dopo}, \emph{during}
\ita{durante}, \emph{until} \ita{fino a}, \emph{without} \ita{senza}, \emph{against}
\ita{contro}, \emph{towards} \ita{verso}.

\needspace{8\baselineskip}\paragraph{Preposizioni articolate.} English has none: the preposition and the article
stay two separate words (two tokens).
\begin{center}
\footnotesize
\begin{tabular}{@{}l l l l l@{}}
\toprule
\textbf{Italiano} & \ita{del} (di + il) & \ita{al} (a + il) & \ita{dal} (da + il) & \ita{nel} (in + il) \\
\textbf{English}  & of the / from the   & to the / at the   & from the / by the   & in the \\
\midrule
\textbf{Italiano} & \ita{sul} (su + il) & \ita{della, alla, dalla, nella, sulla} & \ita{dei, ai, dai, nei, sui} & \ldots \\
\textbf{English}  & on the              & of/to/from/in/on the          & of/to/from/in/on the        & \\
\bottomrule
\end{tabular}
\end{center}
In Penn, all prepositions are IN except \emph{to}, which always gets its own tag TO (also
in \emph{to run} \ita{correre}).

%======================================================================
\subsubsection{Conjunctions}
\notinslides

A \textbf{conjunction} connects words, phrases or clauses. \quad
\textbf{Italiano}: \emph{congiunzione}.

\begin{center}
\footnotesize
\begin{tabularx}{\linewidth}{@{}>{\raggedright\arraybackslash}p{2.6cm} l l >{\raggedright\arraybackslash}X@{}}
\toprule
\textbf{Type} & \textbf{English} & \textbf{Italiano} & \textbf{Example} \\
\midrule
coordinating     & and        & \ita{e}         & \emph{bread and butter} \ita{pane e burro} \\
\emph{(coordinanti)}, CC & or & \ita{o, oppure} & \emph{tea or coffee?} \ita{tè o caffè?} \\
                 & but        & \ita{ma, però}  & \emph{tired but happy} \ita{stanco ma felice} \\
                 & nor        & \ita{né}        & \emph{neither Bob nor Mary} \ita{né Bob né Maria} \\
                 & so         & \ita{quindi, perciò} & \emph{it rained, so I stayed} \ita{pioveva, quindi sono rimasto} \\
\midrule
subordinating    & because    & \ita{perché, poiché} & \emph{I stayed because it rained} \ita{sono rimasto perché pioveva} \\
\emph{(subordinanti)}, IN & although & \ita{sebbene, anche se} & \emph{although it was late} \ita{sebbene fosse tardi} \\
                 & if         & \ita{se}        & \emph{if you want} \ita{se vuoi} \\
                 & when       & \ita{quando}    & \emph{when I arrived} \ita{quando sono arrivato} \\
                 & while      & \ita{mentre}    & \emph{while she slept} \ita{mentre dormiva} \\
                 & that       & \ita{che}       & \emph{I know that you are right} \ita{so che hai ragione} \\
                 & until      & \ita{finché}    & \emph{wait until I call} \ita{aspetta finché non chiamo} \\
                 & so that    & \ita{affinché}  & \emph{so that you understand} \ita{affinché tu capisca} \\
                 & as / since & \ita{siccome}   & \emph{as it was late, \ldots} \ita{siccome era tardi, \ldots} \\
\bottomrule
\end{tabularx}
\end{center}
Penn uses IN for both prepositions and subordinating conjunctions: \emph{before lunch}
\ita{prima di pranzo} and \emph{before I left} \ita{prima che partissi} get the same tag.

%======================================================================
\subsubsection{Interjections, numerals and particles}
\notinslides

\begin{center}
\footnotesize
\begin{tabularx}{\linewidth}{@{}>{\raggedright\arraybackslash}p{3.3cm} >{\raggedright\arraybackslash}X l@{}}
\toprule
\textbf{Type} & \textbf{English $\to$ Italiano} & \textbf{Penn} \\
\midrule
interjection \emph{(interiezione)} & oh \ita{oh}, wow \ita{caspita}, ouch \ita{ahi}, hey \ita{ehi}, yes \ita{sì}, no \ita{no}, hello \ita{ciao} & UH \\
cardinal \emph{(numerale cardinale)} & one \ita{uno}, three \ita{tre}, twenty \ita{venti}, 1908, 3.14 \ita{3,14} & CD \\
ordinal \emph{(numerale ordinale)}  & first \ita{primo}, second \ita{secondo}, third \ita{terzo} & JJ \\
particle (phrasal verbs)         & give \underline{up} \ita{rinunciare}, turn \underline{off} \ita{spegnere}, go \underline{away} \ita{andare via} & RP \\
\bottomrule
\end{tabularx}
\end{center}
Particles look like prepositions but belong to the verb, changing its meaning. No
Italian part of speech corresponds; the closest is an adverb (\ita{andare
\underline{via}}, \ita{tirare \underline{su}}). Note the decimal separator: \emph{3.14} in
English, \ita{3,14} in Italian, which matters for tokenisers
(Section~\ref{sec:sentence-breaking}).

%======================================================================
\subsubsection{Worked example: analisi grammaticale in English}
\notinslides

\begin{nbexample}{tagging a sentence}
\emph{Wow, there are three very big dogs in your garden, but my sister did not see them
because she was running.}\\
\ita{Caspita, ci sono tre cani molto grandi nel tuo giardino, ma mia sorella non li ha
visti perché stava correndo.}
\begin{center}
\footnotesize
\setlength{\tabcolsep}{3pt}
\begin{tabular}{@{}l l l l @{\hspace{0.25cm}} l l l l@{}}
\toprule
\textbf{Word} & \textbf{it.} & \textbf{Tag} & \textbf{Parte del discorso} & \textbf{Word} & \textbf{it.} & \textbf{Tag} & \textbf{Parte del discorso} \\
\midrule
Wow    & \ita{caspita} & UH   & interiezione          & my      & \ita{mia}     & PRP\$ & agg.\ possessivo \\
there  & \ita{ci}      & EX   & (sogg.\ fittizio)     & sister  & \ita{sorella} & NN   & nome comune \\
are    & \ita{sono}    & VBP  & verbo (essere)        & did     & --            & VBD  & verbo ausiliare \\
three  & \ita{tre}     & CD   & agg.\ numerale        & not     & \ita{non}     & RB   & avv.\ di negazione \\
very   & \ita{molto}   & RB   & avv.\ di quantità     & see     & \ita{visti}   & VB   & verbo \\
big    & \ita{grandi}  & JJ   & agg.\ qualificativo   & them    & \ita{li}      & PRP  & pron.\ pers.\ compl. \\
dogs   & \ita{cani}    & NNS  & nome comune, plur.    & because & \ita{perché}  & IN   & cong.\ subordinante \\
in     & \ita{nel}     & IN   & preposizione          & she     & --            & PRP  & pron.\ pers.\ sogg. \\
your   & \ita{tuo}     & PRP\$ & agg.\ possessivo     & was     & \ita{stava}   & VBD  & verbo ausiliare \\
garden & \ita{giardino} & NN  & nome comune           & running & \ita{correndo} & VBG & verbo (gerundio) \\
but    & \ita{ma}      & CC   & cong.\ coordinante    &         &               &      & \\
\bottomrule
\end{tabular}
\end{center}
The ``--'' cells show two differences from Italian: the auxiliary \emph{did} has no
Italian counterpart (Italian negates with \ita{non} alone), and the subject \emph{she} is
dropped in Italian (\ita{stava correndo}). Note also \emph{in your} $\to$ \ita{nel tuo}
(preposizione articolata).

A statistical tagger (NLTK's default one) gets most of this right but tags \emph{Wow} as
NNP (a proper noun): a capitalised unknown word at the start of a sentence looks like a
name, and interjections are rare in the newspaper text taggers are trained on.
\end{nbexample}

%======================================================================
\subsubsection{From words to sentences: analisi logica}
\notinslides

Parts of speech classify \emph{words}. The \textbf{functions} words and phrases play in a
sentence are the subject of \emph{analisi logica}; in the parse trees of
Section~\ref{sec:grammars-parsing} they correspond to positions in the tree. Running
example: \emph{Cynthia sold the bike to Bob for \$200 yesterday} \ita{ieri Cynthia ha
venduto la bici a Bob per 200 dollari}.

\begin{center}
\footnotesize
\begin{tabularx}{\linewidth}{@{}>{\raggedright\arraybackslash}p{2.8cm} >{\raggedright\arraybackslash}p{3.0cm} >{\raggedright\arraybackslash}X@{}}
\toprule
\textbf{English} & \textbf{Italiano} & \textbf{In the parse tree / example} \\
\midrule
Subject           & soggetto            & the $NP$ before the verb: \emph{\underline{Cynthia}} \ita{Cynthia} \\
Predicate (verb)  & predicato verbale   & the verb of the $VP$: \emph{sold} \ita{ha venduto} \\
Predicate (\emph{be} + noun/adj.) & predicato nominale & \emph{she \underline{is happy}} \ita{lei è felice} (\emph{happy} = nome del predicato) \\
Direct object     & complemento oggetto & the $NP$ after the verb: \emph{\underline{the bike}} \ita{la bici} \\
Indirect object   & complemento di termine & \emph{\underline{to Bob}} \ita{a Bob}; also without \emph{to}: \emph{gave \underline{Bob} the bike} \ita{diede la bici a Bob} \\
Adverbials        & complementi indiretti (prezzo, tempo, luogo, \ldots) & $PP$s and adverbs: \emph{\underline{for \$200}} \ita{per 200 dollari}, \emph{\underline{yesterday}} \ita{ieri}, \emph{\underline{in Genoa}} \ita{a Genova} \\
\bottomrule
\end{tabularx}
\end{center}

\begin{nbwhy}
These functions are what semantic role labelling turns into roles
(Section~\ref{sec:frame-semantics}): in the running example the subject is the
\textsc{seller}, the direct object the \textsc{goods}, the indirect object the
\textsc{buyer}, the \emph{for}-phrase (\emph{complemento di prezzo}) the \textsc{price}.
\end{nbwhy}

\begin{nbsummary}[parts of speech]
\begin{itemize}[nosep]
  \item Open classes: \textbf{noun} (NN\ldots{} = nome), \textbf{verb} (VB\ldots, MD = verbo),
        \textbf{adjective} (JJ\ldots{} = agg.\ qualificativo), \textbf{adverb} (RB\ldots{} =
        avverbio).
  \item Closed classes: \textbf{pronoun} (PRP, WP = pronome), \textbf{determiner} (DT =
        articolo, agg.\ determinativi), \textbf{preposition} (IN, TO = preposizione),
        \textbf{conjunction} (CC coordinating, IN subordinating = congiunzione),
        \textbf{interjection} (UH = interiezione); plus numerals (CD) and particles (RP).
  \item English vs.\ Italian: no gender, invariable articles and adjectives, adjectives
        before the noun, subject pronoun always required, no preposizioni articolate,
        prepositions do not map one-to-one, few verb forms (the rest via auxiliaries and
        modals), \emph{do} in questions and negations.
  \item Analisi grammaticale = POS tagging; analisi logica = functions in the sentence
        (subject, predicate, objects, adverbials).
\end{itemize}
\end{nbsummary}

\subsection{Corpora, Genre and Collocations}
\label{sec:corpora}
\slidesrc{Terminology introduction: 39--41}

\subsubsection{Corpus}
\slidesrc{Terminology introduction: 39}

\begin{nbdef}{Corpus (plural: corpora)}
A \textbf{text corpus} is a large and structured set of texts (nowadays usually stored
and processed electronically). Corpora are used for statistical analysis and hypothesis
testing, for checking occurrences, and for validating linguistic rules within a specific
language territory.
\end{nbdef}

Example: the corpora by Mark Davies, originally at \url{https://corpus.byu.edu/} and now
at \url{https://www.english-corpora.org/}: COCA (Corpus of Contemporary American English),
NOW (News on the Web), and others. You can search them to see how often and in which
contexts a word or phrase is used.

\begin{nbwhy}
Corpora are the raw material of statistical and neural NLP
(Section~\ref{sec:nlp-history}): probabilities, word embeddings and LLMs are all learnt
from corpora. What the corpus contains (its size, languages, genres, time period)
determines what the model learns.
\end{nbwhy}

\subsubsection{Genre}
\slidesrc{Terminology introduction: 40}

\begin{nbdef}{Genre}
\begin{itemize}[nosep]
  \item \textbf{Genre}: a style or category of art, music or literature.
  \item \textbf{Genre of a corpus}: usually defined on the basis of
        \emph{text-external} features, such as medium, function and format (e.g.\
        newspaper articles, academic papers, fiction, spoken conversation transcripts,
        social media posts).
\end{itemize}
\end{nbdef}

``Text-external'' means that the genre is decided by \emph{where the text comes from and
what it is for}, not by counting words inside it. Texts of different genres differ a
lot (sentence length, vocabulary, formality), so a model trained on newspapers often
performs worse on tweets.

\subsubsection{Collocations}
\slidesrc{Terminology introduction: 41}

\begin{nbdef}{Collocation}
A \textbf{collocation} is a sequence of words or terms that co-occur \emph{more often
than would be expected by chance}. Examples: ``pay attention'', ``fast food'', ``make an
effort'', ``powerful engine''.
\end{nbdef}

\emph{Fast} and \emph{quick} mean roughly the same thing, yet native speakers combine
them with different nouns (example from the NOW corpus):

\begin{center}
\begin{tabular}{@{}l l@{}}
\toprule
\textbf{Natural English} & \textbf{Unnatural English} \\
\midrule
the fast train  & the \sout{quick} train \\
fast food       & \sout{quick} food \\
a quick shower  & a \sout{fast} shower \\
a quick meal    & a \sout{fast} meal \\
\bottomrule
\end{tabular}
\end{center}

\begin{nbintuition}[``more often than by chance'']
If words were combined at random, \emph{fast food} and \emph{quick food} would appear
about as often as the frequencies of \emph{fast} and \emph{quick} predict. In a corpus,
\emph{fast food} appears far more often than that, and \emph{quick food} hardly ever.
That excess is what makes it a collocation. Nothing in grammar or meaning forbids
``quick food'': collocations are a matter of \emph{usage}, which is why they are learnt
from corpora.
\end{nbintuition}

\subsection{Precision and Recall}
\label{sec:precision-recall}
\slidesrc{Terminology introduction: 43--51; Semantics (long version): 44--50}

Suppose a program searches a collection of documents and returns those it thinks are
relevant to a query. How good is it? Two questions matter, and they pull in opposite
directions:
\begin{itemize}[nosep]
  \item of what it \emph{returned}, how much is actually relevant? (\textbf{precision})
  \item of what is actually relevant, how much did it \emph{find}? (\textbf{recall})
\end{itemize}

\begin{center}
\begin{tikzpicture}
  \draw[thick, nbgray] (-3.8,-2.2) rectangle (4.8,2.3);
  \fill[nbblue!15]   (-0.7,0) circle (1.6);
  \fill[nborange!20] (1.7,0)  circle (1.6);
  \begin{scope}
    \clip (-0.7,0) circle (1.6);
    \fill[nbgreen!35] (1.7,0) circle (1.6);
  \end{scope}
  \draw[thick, nbblue]   (-0.7,0) circle (1.6);
  \draw[thick, nborange] (1.7,0)  circle (1.6);
  \node[nbblue, font=\small\bfseries]   at (-1.6,1.9) {relevant};
  \node[nborange, font=\small\bfseries] at (2.6,1.9)  {retrieved};
  \node[font=\footnotesize, align=center] at (-1.3,0) {relevant,\\missed\\(FN)};
  \node[font=\footnotesize\bfseries] at (0.5,0) {TP};
  \node[font=\footnotesize, align=center] at (2.3,0) {retrieved,\\not\\relevant\\(FP)};
  \node[font=\footnotesize, nbgray, anchor=south west] at (-3.8,-2.2) {corpus (all documents)};
  \node[font=\footnotesize, anchor=south east] at (4.8,-2.2) {neither (TN)};
\end{tikzpicture}
\captionof{figure}{The green overlap (TP) holds the documents that are both relevant
and retrieved. Precision looks at the orange circle (what fraction of it is green),
recall looks at the blue circle (what fraction of it is green).}
\end{center}

\begin{nbdef}{Precision and recall}
\begin{align*}
  \text{precision} &= \frac{\abs{\{\text{relevant documents}\} \cap \{\text{retrieved documents}\}}}
                           {\abs{\{\text{retrieved documents}\}}} \\[4pt]
  \text{recall}    &= \frac{\abs{\{\text{relevant documents}\} \cap \{\text{retrieved documents}\}}}
                           {\abs{\{\text{relevant documents}\}}}
\end{align*}
\textbf{Precision} (also called \emph{positive predictive value}) is the fraction of
retrieved documents that are relevant to the query. \textbf{Recall} (also called
\emph{sensitivity}) is the fraction of relevant documents that are successfully
retrieved. Both are between 0 and 1; $\abs{S}$ denotes the number of elements of the
set $S$.
\end{nbdef}

The numerator is the \emph{same} in both formulas: the documents that are both relevant
and retrieved. Only the denominator changes: precision divides by what the system
\emph{returned}, recall by what it \emph{should have returned}.

\subsubsection{The same thing with a confusion matrix}
\notinslides

Each document falls in one of four cases:

\begin{center}
\begin{tabular}{@{}l c c@{}}
\toprule
 & \textbf{relevant} & \textbf{not relevant} \\
\midrule
\textbf{retrieved}     & true positive (TP)  & false positive (FP) \\
\textbf{not retrieved} & false negative (FN) & true negative (TN) \\
\bottomrule
\end{tabular}
\end{center}

\[
  \text{precision} = \frac{TP}{TP + FP}, \qquad
  \text{recall}    = \frac{TP}{TP + FN}.
\]
Neither formula uses TN: the (usually huge) number of irrelevant documents that were
correctly left out does not count.

\subsubsection{Four worked examples}
\slidesrc{Terminology introduction: 44--51}

The setting is the same for all four examples. You implement an application that
retrieves the documents containing one specific word (``cat'') and its inflectional
forms (``cats'', see Section~\ref{sec:morphology}) in a corpus of \textbf{1000}
documents. Exactly \textbf{24} documents contain ``cat'': these are the \textbf{relevant}
documents.

\begin{nbexample}{1: a typical system}
The application retrieves \textbf{36} documents; \textbf{16} of them contain ``cat'' and
20 do not. So $TP=16$, $FP=20$, $FN = 24-16 = 8$, $TN = 1000-16-20-8 = 956$.
\begin{align*}
  \text{precision} &= \frac{16}{36} = \frac{4}{9} \approx 0.444
     \reason{$TP/(TP+FP) = 16/(16+20)$} \\
  \text{recall}    &= \frac{16}{24} = \frac{2}{3} \approx 0.667
     \reason{$TP/(TP+FN) = 16/(16+8)$}
\end{align*}
Fewer than half of the results are right, and a third of the cat documents are missed.
\end{nbexample}

\begin{nbexample}{2: return everything}
The application retrieves \textbf{all 1000} documents; \textbf{24} of them contain
``cat'' and 976 do not. So $TP = 24$, $FP = 976$, $FN = 0$, $TN = 0$.
\begin{align*}
  \text{precision} &= \frac{24}{1000} = 0.024 \\
  \text{recall}    &= \frac{24}{24} = 1
\end{align*}
Perfect recall, for free: a system that returns everything never misses anything. Its
precision is terrible.
\end{nbexample}

\begin{nbwarning}[a typo in the slides]
Example~2 on the slides says the retrieved documents ``out of which 24 contain the word
cat, and 1986 do not contain it''. Since only 1000 documents are retrieved, the correct
number is $1000 - 24 = 976$. The precision and recall values on the slides are correct.
\end{nbwarning}

\begin{nbexample}{3: play it safe}
The application retrieves only \textbf{2} documents, and both contain ``cat''. So
$TP=2$, $FP=0$, $FN=22$, $TN=976$.
\begin{align*}
  \text{precision} &= \frac{2}{2} = 1 \\
  \text{recall}    &= \frac{2}{24} = \frac{1}{12} \approx 0.083
\end{align*}
Perfect precision, almost for free: return only the one or two results you are
completely sure about. Its recall is terrible.
\end{nbexample}

\begin{nbexample}{4: the perfect system}
The application retrieves \textbf{24} documents, all containing ``cat''. So $TP=24$,
$FP=0$, $FN=0$, $TN=976$.
\[
  \text{precision} = \frac{24}{24} = 1, \qquad \text{recall} = \frac{24}{24} = 1.
\]
\end{nbexample}

\begin{center}
\begin{tikzpicture}
\begin{axis}[ybar, bar width=13pt, width=\linewidth, height=6.5cm,
    ymin=0, ymax=1.15, ytick={0,0.25,0.5,0.75,1},
    symbolic x coords={Ex.\ 1, Ex.\ 2, Ex.\ 3, Ex.\ 4}, xtick=data,
    tick label style={font=\scriptsize}, label style={font=\small},
    legend style={font=\small, at={(0.5,1.03)}, anchor=south, legend columns=3, draw=none},
    enlarge x limits=0.18,
    nodes near coords, every node near coord/.append style={font=\scriptsize, /pgf/number format/fixed, /pgf/number format/precision=2}]
  \addplot[fill=nborange!70, draw=nborange] coordinates {(Ex.\ 1,0.444) (Ex.\ 2,0.024) (Ex.\ 3,1) (Ex.\ 4,1)};
  \addplot[fill=nbblue!60, draw=nbblue]     coordinates {(Ex.\ 1,0.667) (Ex.\ 2,1) (Ex.\ 3,0.083) (Ex.\ 4,1)};
  \addplot[fill=nbgreen!60, draw=nbgreen]   coordinates {(Ex.\ 1,0.533) (Ex.\ 2,0.047) (Ex.\ 3,0.154) (Ex.\ 4,1)};
  \legend{precision, recall, $F_1$}
\end{axis}
\end{tikzpicture}
\captionof{figure}{Precision, recall and $F_1$ for the four examples. Examples 2 and 3
each score perfectly on one measure with a trivial strategy; only $F_1$ exposes them.}
\end{center}

\begin{nbintuition}[the trade-off]
Precision and recall must always be read \emph{together}. Returning more documents
tends to raise recall and lower precision (Example~2 is the extreme); returning fewer
tends to raise precision and lower recall (Example~3). A single number that is high only
when \emph{both} are high is the $F_1$ score below.
\end{nbintuition}

\begin{nbdef}{$F_1$ score (not on the slides)}
The \textbf{$F_1$ score} is the harmonic mean of precision $P$ and recall $R$:
\[
  F_1 = \frac{2PR}{P + R} = \frac{2\,TP}{2\,TP + FP + FN}.
\]
The harmonic mean is dominated by the smaller of the two values, so a system cannot get a
good $F_1$ by maximising only one of them.
\end{nbdef}

For the four examples:
\begin{align*}
  \text{Ex.\ 1:}\quad F_1 &= \frac{2\cdot 16}{2\cdot 16 + 20 + 8} = \frac{32}{60} \approx 0.533 \\
  \text{Ex.\ 2:}\quad F_1 &= \frac{2\cdot 24}{2\cdot 24 + 976 + 0} = \frac{48}{1024} \approx 0.047 \\
  \text{Ex.\ 3:}\quad F_1 &= \frac{2\cdot 2}{2\cdot 2 + 0 + 22} = \frac{4}{26} \approx 0.154 \\
  \text{Ex.\ 4:}\quad F_1 &= \frac{2\cdot 24}{2\cdot 24 + 0 + 0} = 1
\end{align*}

\subsubsection{Another example: five algorithms}
\slidesrc{Semantics (long version): 44--50}

The long version of the semantics slides repeats the definitions with a second example.
A query asks for the documents about topic $T$ in a corpus $C$. The correct answer is
known: $D_1, D_2, D_4, D_5, D_6, D_7, D_8$ deal with $T$ (7 relevant documents) and no
other document does.

\begin{center}
\footnotesize
\begin{tabular}{@{}l l c c c@{}}
\toprule
\textbf{Algorithm} & \textbf{Returns} & $\abs{\text{rel}\cap\text{ret}}$ & \textbf{Precision} & \textbf{Recall} \\
\midrule
A1 & $D_1, D_2, D_3, D_4$            & 3 & $3/4 = 75\%$       & $3/7 \approx 43\%$ \\
A2 & $D_1, \dots, D_8$               & 7 & $7/8 = 87.5\%$     & $7/7 = 100\%$ \\
A3 & $D_1$ only                      & 1 & $1/1 = 100\%$      & $1/7 \approx 14\%$ \\
A4 & all and only the relevant & 7 & $7/7 = 100\%$      & $7/7 = 100\%$ \\
A5 & 1000 docs, incl.\ the 7 relevant & 7 & $7/1000 = 0.7\%$ & $7/7 = 100\%$ \\
\bottomrule
\end{tabular}
\end{center}

The same lessons as before: A3 buys perfect precision by returning almost nothing, A5
buys perfect recall by returning almost everything; only A4 is perfect on both. (On the
slide for A2 the formulas are labelled ``precision(A1)'' and ``recall(A1)'': a
copy-paste slip, the values are those of A2.)

\subsubsection{Exercises}
\notinslides

\begin{nbexercise}{1}{morphemes}
Split each word into morphemes and classify each morpheme as free or bound, and each
bound morpheme as derivational or inflectional: \emph{unhappiness}, \emph{walked},
\emph{teachers}, \emph{rewritten}. Which of \emph{bookshelf}, \emph{bookish},
\emph{books} is a compound?
\end{nbexercise}
\begin{nbsolution}
\begin{itemize}[nosep]
  \item \emph{un-happi-ness}: \emph{un-} bound, derivational (negates the meaning);
        \emph{happy} free root; \emph{-ness} bound, derivational (adjective $\to$ noun).
  \item \emph{walk-ed}: \emph{walk} free root; \emph{-ed} bound, inflectional (past tense).
  \item \emph{teach-er-s}: \emph{teach} free root; \emph{-er} bound, derivational
        (verb $\to$ noun, ``one who teaches''); \emph{-s} bound, inflectional (plural).
  \item \emph{re-writ-ten}: \emph{re-} bound, derivational (``again''); \emph{write}
        free root; \emph{-en} bound, inflectional (past participle).
\end{itemize}
Only \emph{bookshelf} is a compound (\emph{book} + \emph{shelf}, both free).
\emph{Bookish} is derived (\emph{-ish} is bound), \emph{books} is inflected.

\emph{Lesson}: derivational morphemes sit closer to the root than inflectional ones
(\emph{teach-er-s}, not \emph{teach-s-er}).
\end{nbsolution}

\begin{nbexercise}{2}{precision and recall}
A corpus has 500 documents, 40 of which are relevant to a query. A search engine
returns 50 documents, 30 of which are relevant. Compute TP, FP, FN, TN, precision,
recall and $F_1$. What happens to precision and recall if the engine also returns 50 more
documents, none of them relevant?
\end{nbexercise}
\begin{nbsolution}
$TP = 30$, $FP = 50 - 30 = 20$, $FN = 40 - 30 = 10$, $TN = 500 - 30 - 20 - 10 = 440$.
\begin{align*}
  \text{precision} &= \frac{30}{50} = 0.6, \qquad
  \text{recall} = \frac{30}{40} = 0.75, \\
  F_1 &= \frac{2\cdot 30}{2\cdot 30 + 20 + 10} = \frac{60}{90} \approx 0.667.
\end{align*}
With 50 more irrelevant documents: $TP = 30$, $FP = 70$, $FN = 10$, so
precision $= 30/100 = 0.3$ (halved) while recall stays $30/40 = 0.75$, and
$F_1 = 60/(60+70+10) = 60/140 \approx 0.429$.

\emph{Lesson}: adding wrong results never changes recall but always lowers precision.
\end{nbsolution}

\begin{nbsummary}[terminology]
\begin{itemize}[nosep]
  \item \textbf{Morpheme}: smallest meaningful unit; \textbf{free} (stands alone:
        \emph{dog}) or \textbf{bound} (\emph{un-}, \emph{-ness}, \emph{-s}).
  \item Bound morphemes: \textbf{derivational} (new word: meaning or part of speech
        changes, \emph{happy} $\to$ \emph{happiness}) or \textbf{inflectional} (new form
        of the same word: tense, number, \ldots; \emph{dog} $\to$ \emph{dogs}).
  \item \textbf{Word}: smallest element uttered in isolation with meaning; one or more
        morphemes. \textbf{Compound}: word made of $\geq 2$ \emph{free} morphemes.
  \item \textbf{Corpus}: large structured set of texts. \textbf{Genre of a corpus}: from
        text-external features (medium, function, format). \textbf{Collocation}: words
        co-occurring more often than by chance (\emph{fast food}, not \emph{quick food}).
  \item $\text{precision} = TP/(TP+FP)$: of what I returned, how much is right;
        $\text{recall} = TP/(TP+FN)$: of what is right, how much I returned;
        $F_1 = 2PR/(P+R)$. Returning everything gives recall 1; returning very few
        sure hits gives precision 1.
\end{itemize}
\end{nbsummary}

